% Generated draft: requires mathematical and source review.
\documentclass[11pt]{article}
\usepackage[margin=25mm]{geometry}
\usepackage{fontspec}
\setmainfont{Latin Modern Roman}
\usepackage{amsmath,amssymb,mathtools,unicode-math}
\setmathfont{Latin Modern Math}
\usepackage{enumitem,needspace,xurl,hyperref}
\setlength{\parindent}{0pt}
\setlength{\parskip}{5pt}
\setlength{\emergencystretch}{3em}
\newcommand{\N}{\mathbb{N}}
\newcommand{\Z}{\mathbb{Z}}
\newcommand{\Q}{\mathbb{Q}}
\newcommand{\R}{\mathbb{R}}
\newcommand{\C}{\mathbb{C}}
\newcommand{\F}{\mathbb{F}}
\newcommand{\A}{\mathbb{A}}
\newcommand{\PP}{\mathbb{P}}
\newcommand{\E}{\mathbb{E}}
\newcommand{\eps}{\varepsilon}
\newcommand{\dd}{\,\mathrm{d}}
\begin{document}
\section*{UnsolvedMath ID: 30003771}
\textbf{Source problem code:} catalogue entry
\subsection{Definitions and mathematical statement}
Let \(k\) denote the base field mentioned in the source; its characteristic and further hypotheses are unspecified. A finite \(k\)-group scheme is a group scheme \(G\) whose structure morphism to \(\operatorname{Spec} k\) is finite. It is abelian if its multiplication is commutative. The question concerns allowing finite \(k\)-group schemes without this commutativity assumption.

For a \(k\)-scheme \(U\), put \(G_U=G\times_k U\). A right \(G_U\)-torsor is a \(U\)-scheme \(P\) with a right \(G_U\)-action that is locally isomorphic to \(G_U\), with its translation action, for the faithfully flat and locally finitely presented topology. Ramification requires additional data: an ambient space containing \(U\), a boundary, and a specified notion of extension across that boundary. None of these data, or the meaning of prescribed ramification, is supplied here.

A parabolic vector bundle carries additional filtration and weight data along a specified boundary. The permissible weights and the geometric setting are part of the definition and are not specified in the record. In a neutral Tannakian setting with a chosen fiber functor, the monodromy group scheme of an object is the Tannaka group of the tensor subcategory it generates. In that setting, an object is essentially finite when this monodromy group scheme is finite; abelian monodromy additionally means that this group scheme is commutative. This description is conditional on a suitable Tannakian setting and does not supply the missing hypotheses for the intended parabolic category.

\textbf{Question as supplied.} Does the cited correspondence between tamely ramified torsors under finite abelian \(k\)-group schemes, with prescribed ramification data, and essentially finite parabolic bundles with abelian monodromy admit an analogue for arbitrary finite \(k\)-group schemes \(G\), including nonabelian ones?

The source does not specify the categories, functors, geometric hypotheses, or ramification conditions of the proposed analogue. In particular, it does not say whether the parabolic side should consist of individual objects with finite monodromy, objects equipped with additional \(G\)-structure, or tensor functors from representations of \(G\). Nor does it explicitly state how the abelian-monodromy condition should change. Consequently, the question above preserves the supplied target, but a precise categorical equivalence and its full quantifiers remain unresolved.
\subsection{Short English statement}
The supplied question asks whether the cited abelian torsor--parabolic correspondence extends to finite group schemes that may be noncommutative, while retaining prescribed tame ramification. The geometric hypotheses, ramification data, and exact nonabelian counterpart on the parabolic side are missing, so a fully specified equivalence cannot be stated from this record alone.
\subsection{Sources}
\begin{itemize}
\item Local UnsolvedMath registry: permanent ID \(30003771\), source problem code \texttt{catalogue entry}, title \emph{Nonabelian Torsors and Essentially Finite Parabolic Bundles}, in the Oberwolfach Reports Open Problems collection. The record gives proposed year 2018 and names no individual proposer.
\item Supplied report citation: \emph{Mini-Workshop: Arithmetic Geometry and Symmetries around Galois and Fundamental Groups (2019)}, Oberwolfach Reports, pp.~1021--1080, DOI \url{https://doi.org/10.4171/owr/2018/17}. The supplied citation says 2019, whereas the proposed year and DOI identify 2018; this discrepancy is unresolved.
\item Supplied literature reference: Biswas--Borne, \emph{tamely ramified torsors}, \url{https://arxiv.org/abs/1706.06733}.
\item Supplied literature reference: \emph{Published torsor--parabolic equivalence}, \url{https://doi.org/10.2422/2036-2145.202005_015}.
\item Supplied literature reference: \emph{OWR 2018 report}, \url{https://doi.org/10.4171/owr/2018/17}.
\item Literature assessment dated exactly \texttt{2026-08-21}: the supplied curation attributes the finite abelian equivalence to Biswas--Borne, including a published 2022 version, and labels the extension to arbitrary finite nonabelian group schemes open. This is supplied dated curation derived from the OWR record, not an independently verified literature assessment. The record also supplies the labels \texttt{open} and \texttt{corrected\_verified}; neither is independently established by this draft.
\end{itemize}
\end{document}
